% !TeX root = Lec03_RBC_Basic.tex

\ifdefined\AdvMacroMain
  \chapter{Lecture 3: Real Business Cycle Model}
  \let\ThisNoteEnd\relax
\else
  \documentclass[12pt]{article}
  \newcommand{\ProjectRoot}{..}
  \usepackage[a4paper,margin=0.85in]{geometry}
\usepackage{newpxtext}
\usepackage{amsmath,amssymb,amsthm,mathtools}
\usepackage{newpxmath}
\usepackage{xcolor}
\usepackage{enumitem}
\usepackage{graphicx}
\usepackage{array,tabularx}
\usepackage{float}
\usepackage{hyperref}

\definecolor{noteRed}{RGB}{190,45,90}

\newcommand{\red}[1]{\textcolor{noteRed}{#1}}
\newcommand{\blue}[1]{\textcolor{blue}{#1}}
\newcommand{\purple}[1]{\textcolor{purple}{#1}}

\newcommand{\E}{\mathbb{E}}
\newcommand{\dd}{\mathrm{d}}
\newcommand{\MPN}{\mathrm{MPN}}
\newcommand{\MRS}{\mathrm{MRS}}
\newcommand{\MRT}{\mathrm{MRT}}
\newcommand{\Var}{\operatorname{Var}}
\newcommand{\boxit}[1]{\fbox{\ensuremath{#1}}}
\newcommand{\circled}[1]{\textcircled{\scriptsize #1}}

\newenvironment{tightalign}
  {\[\begin{aligned}}
  {\end{aligned}\]}

\graphicspath{
  {\ProjectRoot/_assets/figures/}
  {\ProjectRoot/_assets/figures/lecture_04/}
}

  \title{Lecture 3: Real Business Cycle Model}
  \author{Chengyuan HE}
  \date{March 2026}
  \begin{document}
  \maketitle
  \def\ThisNoteEnd{\end{document}}
\fi

\subsection*{1. Basic RBC model: Social Planner's Problem}

\textbf{Preference:}
\[
U
=
E_0\!\left[
\sum_{t=0}^{\infty}\beta^t
\frac{C_t^{1-\sigma}-1}{1-\sigma}
\right]
\]
\[
\textcolor{red}{\text{expectation given information at } t=0}
\]

\textbf{Technology:}
\[
C_t + K_t = Z_t \boxed{K_{t-1}}^{\alpha} N_t^{1-\alpha} + (1-\delta)K_{t-1}
\]
\[
\textcolor{red}{ \text{The amount of capital is predetermined in last period.}}
\]
\[
\textcolor{red}{\text{Produced goods \& depreciated capital in last period used for consumption \& new capital.}}
\]
\[
\textcolor{red}{\text{In current period } t,\ \text{agent decides on } K_t.}
\]
\[
Z_t \text{ is the TFP which follows the process:}
\]
\[
\log Z_t
=
(1-\psi)\underbrace{\log \bar Z}_{\textcolor{red}{\displaystyle \text{steady state TFP}}}
+
\psi \log Z_{t-1}
+
\varepsilon_t
\]
\[
\textcolor{red}{\text{usually normalize } \bar Z=1 \Rightarrow \log \bar Z=0}
\]

\textbf{Endowment:}
\[
N_t=1,\qquad K_{-1}>0
\]
\[
\textcolor{red}{\text{time/labor endowment}\qquad \text{initial capital stock}}
\]

\textbf{Information:} decision made based on all information $I_t$ up to $t$.

Social planner maximize the social utility by choosing $C_t,K_t$:
\[
\max_{\{C_t,K_t\}} U
=
\max_{\{C_t,K_t\}}
E_0\!\left[
\sum_{t=0}^{\infty}\beta^t
\frac{C_t^{1-\sigma}-1}{1-\sigma}
\right]
\]

subject to a budget constraint:
\[
\forall t,\qquad
C_t+K_t
=
Z_tK_{t-1}^{\alpha}N_t^{1-\alpha}
+
(1-\delta)K_{t-1}
\]
\[
N_t=1,\qquad K_{-1}>0
\]

Set up dynamic Lagrangian equation:
\[
\mathcal L
=
E_0\sum_{t=0}^{\infty}\beta^t
\left\{
\frac{C_t^{1-\sigma}-1}{1-\sigma}
+
\lambda_t\bigl[
Z_tK_{t-1}^{\alpha}
+
(1-\delta)K_{t-1}
-
C_t-K_t
\bigr]
\right\}
\]

F.O.C.:
\[
\frac{\partial \mathcal L}{\partial C_t}=0
\qquad \Rightarrow \qquad
C_t^{-\sigma}-\lambda_t=0
\tag{1}
\]
\[
\begin{aligned}
\frac{\partial \mathcal L}{\partial K_t}=0
\qquad \Rightarrow \qquad
&
-\lambda_t \beta^t
+
E_t\!\left[
\lambda_{t+1}\beta^{t+1}
\bigl(\alpha Z_{t+1}K_t^{\alpha-1}+1-\delta\bigr)
\right]
=0
\\[0.6em]
\Rightarrow\qquad
&
-\lambda_t
+
\beta E_t\!\left[
\lambda_{t+1}
\bigl(\alpha Z_{t+1}K_t^{\alpha-1}+1-\delta\bigr)
\right]
=0
\end{aligned}
\tag{2}
\]
\[
\frac{\partial \mathcal L}{\partial \lambda_t}=0
\qquad \Rightarrow \qquad
C_t+K_t-Z_tK_{t-1}^{\alpha}-(1-\delta)K_{t-1}=0
\]

Transversality Condition:
\[
\lim_{t\to\infty}\beta^t E_t[\lambda_t K_t]=0
\]

Plug in (1) into (2):
\[
\beta E_t\!\left[
C_{t+1}^{-\sigma}
\bigl(\alpha Z_{t+1}K_t^{\alpha-1}+1-\delta\bigr)
\right]
=
C_t^{-\sigma}
\]

\[
\Rightarrow\qquad
1
=
\beta E_t\!\left[
\left(\frac{C_t}{C_{t+1}}\right)^{\sigma}
\bigl(\alpha Z_{t+1}K_t^{\alpha-1}+1-\delta\bigr)
\right]
\]

\newpage

We have a dynamic system of $(K_t,C_t,Z_{t+1})$
\[
\begin{aligned}
K_t
&= Z_t K_{t-1}^{\alpha} + (1-\delta)K_{t-1} - C_t,
&& \text{(1)} \\[0.8em]
1
&= \beta E_t\!\left[\left(\dfrac{C_t}{C_{t+1}}\right)^{\sigma}
\bigl(1-\delta+\alpha Z_{t+1}K_t^{\alpha-1}\bigr)\right],
&& \text{(2)} \\[0.8em]
\log Z_t
&= (1-\psi)\log \bar Z + \psi \log Z_{t-1} + \varepsilon_t,
\qquad \varepsilon_t \sim iid\,N(0,\sigma_\varepsilon^2),
&& \text{(3)}
\end{aligned}
\]
\[
\text{state variables } (K_{t-1},\, C_t),
\qquad
\text{shock variable } Z_t.
\]

Three variables, three equations: \quad \#(equation) = \#(variable)

\subsection*{2. What is the economic intuition behind Lagrangian multiplier?}

Lagrangian multiplier is the shadow price of constrained resources. Marginal value of relaxing or tightening a constraint. For a problem,
\[
V(b)=\max_x f(x)\qquad \text{subject to } g(x)\le b
\]

\[
\mathcal L(x,\lambda)=f(x)+\lambda\,(b-g(x))
\]

\begin{itemize}
    \item If constraint is binding $g(x)=b$, $\lambda$ shows how much the objective function would increase if you relax the constraint slightly.\textcolor{red}{\text{——shadow value}}
\end{itemize}

Envelope theorem:
\[
\frac{dV(b)}{db}
=
\frac{\partial \mathcal L(x,\lambda;b)}{\partial b}\Big|_{x^*,\lambda^*}
\qquad \Rightarrow \qquad
V'(b)=\lambda^*(b)
\]

\[
V(b+\Delta)-V(b)\approx \lambda^*(b)\Delta
\]

\begin{itemize}
    \item If constraint is non-binding $g(x)<b$, $\textcolor{red}{\lambda^*}=0
    \qquad \Rightarrow \qquad V'(b)=0$:
    A relaxation of constraint will not improve the objective function.

\end{itemize}


eg: RBC:
\[
W_t
=
\max_{\{C_{t+j},K_{t+j+1}\}_{j=0}^{\infty}}
E_t\sum_{j=0}^{\infty}\beta^j
\frac{C_{t+j}^{1-\sigma}}{1-\sigma}
\]

subject to
\[
C_{t+j}+K_{t+j+1}
=
F(Z_{t+j},K_{t+j})
+
(1-\delta)K_{t+j}
+\textcolor{red}{a_t}
\]
\[
\textcolor{red}{a_t \text{ appear only at } j=0,\ \text{perturb only one period's constraint}}
\]

\[
\mathcal L
=
E_t\sum_{j=0}^{\infty}\beta^j
\left[
\frac{C_{t+j}^{1-\sigma}}{1-\sigma}
+\lambda_{t+j}
\bigl(F(\cdot)+(1-\delta)K_{t+j}-C_{t+j}-K_{t+j+1}+a_{t+j}\bigr)
\right]
\]

\[
\frac{\partial W_t}{\partial a_t}
=
\frac{\partial \mathcal L_t}{\partial a_t}\Big|_{\text{at optimum}}
=
\lambda_t
\]

\[
\frac{\partial \mathcal L}{\partial C_t}
=
\beta^t C_t^{-\sigma}-\beta^t\lambda_t=0
\qquad \Rightarrow \qquad
C_t^{-\sigma}=\lambda_t
\]

\[
\Rightarrow\qquad
\frac{\partial W_t}{\partial a_t}
=
\lambda_t
=
C_t^{-\sigma}
=
U_C(C_t)
=
\frac{\partial (\text{max attainable lifetime utility } W_t)}{\partial (\text{one more unit of resources at } t)}
\]

\textbf{What is transversality condition (TVC)?}

Boundary condition at infinity that rules out Ponzi-scheme: leaving valuable resources unspent forever.

In a typical RBC model,
\[
\text{TVC:}\qquad
\lim_{T\to\infty}\beta^T U_C(C_T)K_{T+1}=0
\]

The discounted marginal utility value of the capital stock carried into the far future must vanish. Otherwise you're wasting utility by not consuming some of it earlier.

At large $T$, if we increase $C_T$ by $\varepsilon$ and decrease future capital by $\varepsilon$:
\[
\text{First-order loss in capital:}\qquad \beta^{T+1}\lambda_{T+1}\varepsilon
\]
\[
\text{First-order gain in utility:}\qquad \beta^T U_C\varepsilon
\]

\[
\Rightarrow\qquad
\lim_{T\to\infty}\beta^T U_C K_{T+1}=0
\]
\[
\textcolor{red}{\text{in far future, you can not gain by this operation}}
\]

\subsection*{3. Basic RBC model: competitive market approach}

\[
\textcolor{red}{\text{divide social planner into household and firm}}
\]

\textcolor{red}{*Household's problem:}
Household's own capital $K_{t-1}$ \& one unit of labor $N_t=1$. They lend capital and labor to firms and earn interest $R_t$ \& wage $W_t$.

\[
\max_{\{C_t,K_t,N_t\}_{t=0}^{\infty}}
E_0\left[
\sum_{t=0}^{\infty}\beta^t\frac{C_t^{1-\sigma}-1}{1-\sigma}
\right]
\]

s.t.
\[
C_t+K_t^{(s)}
=
W_tN_t^{(s)}
+
R_tK_{t-1}^{(s)}
+
(1-\delta)K_{t-1}^{(s)}
\]
\[
\textcolor{red}{(s)\text{ means supply}}
\]

No-ponzi Scheme condition (TVC):
\[
\lim_{t\to\infty}
E_0\!\left[
\left(\prod_{s=0}^{t}R_s^{-1}\right)K_t^{(s)}
\right]
=0
\]

\textcolor{red}{*Firm's problem:}
Representative firm with tech $Z_t$, hire capital \& labor to produce, and sells output to maximize profit. Given $W_t,R_t$, firm chooses $N_t^{(d)}$ \& $K_{t-1}^{(d)}$.

\[
\Pi_t
=
\max_{\{K_{t-1}^{(d)},N_t^{(d)}\}}
Z_t(K_{t-1}^{(d)})^{\alpha}(N_t^{(d)})^{1-\alpha}
-
W_tN_t^{(d)}
-
R_tK_{t-1}^{(d)}
\]
\[
\textcolor{red}{(d)\text{ means demand}}
\]

where tech follows AR(1):
\[
\log Z_t=(1-\psi)\log \bar Z+\psi \log Z_{t-1}+\varepsilon_t,
\qquad
\varepsilon_t\sim iid\,N(0,\sigma_{\varepsilon}^2)
\]

\textcolor{red}{Solve household's problem with Lagrangian Multiplier Method:}

\[
\mathcal L
=
E_0\sum_{t=0}^{\infty}\beta^t
\left[
\frac{C_t^{1-\sigma}-1}{1-\sigma}
+
\lambda_t\bigl(
R_tK_{t-1}^{(s)}+W_tN_t^{(s)}+(1-\delta)K_{t-1}^{(s)}-C_t-K_t^{(s)}
\bigr)
\right]
\]

\[
\frac{\partial \mathcal L}{\partial C_t}=0
\qquad \Rightarrow \qquad
C_t^{-\sigma}-\lambda_t=0
\tag{1}
\]

\[
\frac{\partial \mathcal L}{\partial K_t}=0
\qquad \Rightarrow \qquad
-\beta^t\lambda_t+\beta^{t+1}E_t[\lambda_{t+1}(R_{t+1}+1-\delta)]=0
\]
\[
\Rightarrow\qquad
\lambda_t
=
\beta E_t[\lambda_{t+1}(R_{t+1}+1-\delta)]
\tag{2}
\]

\[
\frac{\partial \mathcal L}{\partial \lambda_t}=0
\qquad \Rightarrow \qquad
C_t+K_t^{(s)}
=
W_tN_t^{(s)}+R_tK_{t-1}^{(s)}+(1-\delta)K_{t-1}^{(s)}
\tag{3}
\]

Euler equation. Plug (1) into (2):
\[
\underbrace{C_t^{-\sigma}}_{\text{MU of 1 unit } C \text{ today}}
=
\underbrace{\beta}_{\textcolor{red}{\text{discount to today}}}
E_t\!\left[
\underbrace{C_{t+1}^{-\sigma}}_{\text{MU of 1 unit } C \text{ tomorrow}}
\underbrace{(R_{t+1}+1-\delta)}_{\textcolor{red}{\text{return of extra unit of } K \text{ tomorrow}}}
\right]
\]

LHS: The utility of 1 unit of consumption today.

RHS: The expected utility of sacrificing 1 unit of consumption today, instead save as capital, and carry over to next period, get $R_{t+1}$ (rent earning) \& get $(1-\delta)$ unit of depreciated capital, consume $R_{t+1}+1-\delta$ in the next period at marginal utility of consumption $\lambda_{t+1}=C_{t+1}^{-\sigma}$. Finally, discount to current value by multiplying $\beta$.

\[
\text{LHS}=\text{RHS}
\qquad \Rightarrow \qquad
\text{Intertemporal trade-off (``no arbitrage'')}
\]

\begin{itemize}
    \item A TFP shock on $Z_t$ change the level through resource constraint, and dynamics through the Euler equation.
\end{itemize}

Resource constraint:
\[
C_t+K_t=Z_tK_{t-1}^{\alpha}N_t^{1-\alpha}+(1-\delta)K_{t-1}
\]

\textcolor{red}{Solve firm's problem with Lagrangian multiplier method:}

\[
\frac{\partial}{\partial K_{t-1}^{(d)}}
\Bigl[
Z_t(K_{t-1}^{(d)})^{\alpha}(N_t^{(d)})^{1-\alpha}
-W_tN_t^{(d)}-R_tK_{t-1}^{(d)}
\Bigr]
=0
\]

\[
\Rightarrow\qquad
\alpha Z_t(K_{t-1}^{(d)})^{\alpha-1}(N_t^{(d)})^{1-\alpha}=R_t
\tag{4}
\]
\[
\text{MPK}=R_t
\]

\[
\frac{\partial \Pi_t}{\partial N_t^{(d)}}=0
\qquad \Rightarrow \qquad
(1-\alpha)Z_t(K_{t-1}^{(d)})^{\alpha}(N_t^{(d)})^{-\alpha}=W_t
\tag{5}
\]
\[
\text{MPL}=W_t
\]

\textcolor{red}{Close the model: Market clearing conditions}

Capital market:
\[
K_{t-1}^{(d)}=K_{t-1}^{(s)}
\qquad \cdots \quad \text{capital demand = supply}
\]

Labor market:
\[
N_t^{(d)}=N_t^{(s)}=1
\qquad \cdots \quad \text{labor demand = supply}
\]

Goods market:
\[
C_t+K_t=Z_tK_{t-1}^{\alpha}N_t^{1-\alpha}+(1-\delta)K_{t-1}
\qquad \cdots \quad \text{Resource constraint}
\]

Walras Law: if there are $m$ markets, then we need $m-1$ number of market clearing conditions. The last market is automatically cleared.

Thus, here we need two out of three. We keep resource constraint \& capital market clearing by dropping (d) \& (s). We also plug in constant labor supply, demand.

\textcolor{red}{How resource constraint is derived?}

Competitive market $\Rightarrow \Pi_t=0$:
\[
Z_t(K_{t-1}^{(d)})^\alpha (N_t^{(d)})^{1-\alpha}
=
W_tN_t^{(d)}+R_tK_{t-1}^{(d)}
\]

(3) from household's problem:
\[
C_t+K_t^{(s)}
=
W_tN_t^{(s)}+R_tK_{t-1}^{(s)}+(1-\delta)K_{t-1}^{(s)}
\]

\[
(s)=(d)
\Rightarrow
C_t+K_t
=
Z_tK_{t-1}^{\alpha}N_t^{1-\alpha}+(1-\delta)K_{t-1}
\]
\[
\textcolor{red}{\text{Produced goods + depreciated capital from last period is divided into use of capital}}
\]
\[
\textcolor{red}{\text{and consumption}}
\]

\textcolor{red}{System of equations of the competitive market version of the model:}
\[
\text{(i)}\quad
K_t=Z_tK_{t-1}^{\alpha}+(1-\delta)K_{t-1}-C_t
\qquad \cdots \quad \text{resource constraint}
\]

\[
\text{(ii)}\quad
1
=
E_t\left\{
\beta\left(\frac{C_t}{C_{t+1}}\right)^{\sigma}
\left[
\alpha Z_{t+1}K_t^{\alpha-1}+(1-\delta)
\right]
\right\}
\]

\[
\text{(iii)}\quad
\log Z_t=(1-\psi)\log \bar Z+\psi \log Z_{t-1}+\varepsilon_t,\qquad
\varepsilon_t\sim iid\,N(0,\sigma_{\varepsilon}^2)
\]

\[
\textcolor{red}{\ast\ \text{The competitive market version is equivalent to Social planner's version. (isomorphic)}}
\]
\[
\textcolor{red}{\Rightarrow\ \text{Same equilibrium outcome}}
\]

\textcolor{red}{* The competitive market version of RBC model is a decentralized version of social planner RBC.}
\newpage
\subsection*{4. First and second welfare theorem in dynamic equilibrium}

*First welfare theorem. Provided that all product and factor markets are \textcolor{red}{competitive} and there are \textcolor{red}{no externalities} in production or consumption, dynamic competitive equilibrium are Pareto Optimal.
\[
\textcolor{red}{\text{``no friction'' world in economics}}
\]

*Second welfare theorem. All Pareto optimal allocation can be decentralized as a dynamic competitive equilibrium.

\noindent $1^{st}$ Welfare theorem in RBC:\quad
CE allocation solves the social planner's problem
$\Rightarrow$
Competitive equilibria are Pareto Optimal.

\medskip

\noindent $2^{nd}$ Welfare theorem in RBC:\quad
Social planner allocation solves the competitive equilibrium
$\Rightarrow$
Pareto efficient allocation is supported by CE ($\exists$ competitive prices that support it).


Both theorems rely on:
\begin{enumerate}[label=\arabic*)]
    \item perfect competition (price takers)
    \item No externalities, no taxes/distortions, no market power.
    \item Complete contracting:no missing markets (rep agent RBC avoids this)
    \item Concave utility, convex technology set (CRS + concavity)
    \item TVC / no-Ponzi to rule out pathological asset plans
\end{enumerate}
\newpage
\textcolor{red}{Question 1. If there is a exogenous TFP shock $Z_t$, which moves more on impact $C_t$ or $K_t$}


\[
Y_t = C_t + i_t
\qquad \Rightarrow \qquad
\Delta Y = \Delta C + \Delta i
\]

\[
Y_t = Z_t K_{t-1}^{\alpha}
\qquad \Rightarrow \qquad
Z_t \uparrow \Rightarrow Y_t \uparrow
\]

\[
\log Z_{t+1} = (1-\psi)\log \bar Z + \psi \log Z_t + \varepsilon_t
\qquad \Rightarrow \qquad
Z_t \uparrow \Rightarrow Z_{t+1} \uparrow \Rightarrow MPK_{t+1} \uparrow \Rightarrow R_{t+1} \uparrow
\]

\[
i_t = K_t - (1-\delta)K_{t-1}
\qquad \Rightarrow \qquad
K_t \uparrow \Rightarrow i_t \uparrow
\]

\[
1 = \beta E_t \left[
\left(\frac{C_t}{C_{t+1}}\right)^{\sigma}
(R_{t+1}+1-\delta)
\right]
\qquad \Rightarrow \qquad
R_{t+1} \uparrow \Rightarrow \frac{C_t}{C_{t+1}} \downarrow \Rightarrow C_{t+1} \uparrow > C_t \uparrow
\]

Thus, \qquad
$Z_t \uparrow \Rightarrow C_t \uparrow,\ C_{t+1} \uparrow \ \&\ i_t \uparrow$

$K_{t-1}$ is pre-determined, thus at current period $t$, $\Delta K_{t-1}=0$.

\[
i_t \uparrow \gg C_t \uparrow \; ? \qquad \text{Ambiguous for now.}
\]

1) but in data, it is true

2) 1$^\circ$\quad $\sigma$ captures intertemporal substitution

\[
\sigma \uparrow
\qquad \Rightarrow \qquad
\text{stronger consumption smoothing}
\qquad \Rightarrow \qquad
i_t \uparrow \gg C_t \uparrow
\]
\[
\text{more likely to hold}
\]

2$^\circ$ \quad $\psi$ captures shock persistence

\[
\psi \uparrow
\qquad \Rightarrow \qquad
i_t \uparrow \gg C_t \uparrow
\qquad \text{more likely to hold}
\]

3$^\circ$ \quad $\alpha$ captures MPK

\[
\alpha \uparrow
\qquad \Rightarrow \qquad
i_t \uparrow \gg C_t \uparrow
\qquad \text{more likely}
\]

4$^\circ$ \quad $\delta$ captures depreciation

\[
\delta \uparrow
\qquad \Rightarrow \qquad
i_t \uparrow \gg C_t \uparrow
\qquad \text{less likely}
\]

5$^\circ$ \quad $\beta$ captures patience

\[
\beta \uparrow
\qquad \Rightarrow \qquad
i_t \uparrow \gg C_t \uparrow
\qquad \text{more likely}
\]

But in normal calibration, where
\[
\sigma = 2,\qquad
\psi = 0.95,\qquad
\alpha = \frac{1}{3},
\]
\[
\delta = 0.025 \; (\text{quarter data}),\qquad
\beta = 0.99 \; (\text{quarter data})
\]

\[
i_t \uparrow \gg C_t \uparrow
\qquad \text{and} \qquad
std(i_t) > std(C_t)
\]

\textcolor{red}{Question 2: What if households sell capital to firms and firms sell back to household after production?}

\textbf{Household's problem:}
\[
\max_{\{C_t,K_t,N_t\}_{t=0}^{\infty}}
E_0\left[
\sum_{t=0}^{\infty}\beta^t\frac{C_t^{1-\sigma}-1}{1-\sigma}
\right]
\]

s.t.
\[
C_t+K_t^{(s)}=W_tN_t^{(s)}+R_tK_{t-1}^{(s)}
\]

\[
\Rightarrow\qquad
\begin{aligned}
C_t^{-\sigma}
&=\beta E_t\left[C_{t+1}^{-\sigma}R_{t+1}\right],
&& \text{(1)}\\[0.5em]
C_t+K_t^{(s)}
&=W_tN_t^{(s)}+R_tK_{t-1}^{(s)},
&& \text{(2)}
\end{aligned}
\]

\textbf{Firm's problem:}
\[
\max_{\{K_{t-1}^{(d)},N_t^{(d)}\}}
\Pi_t
=
\max
\left\{
Z_t(K_{t-1}^{(d)})^\alpha (N_t^{(d)})^{1-\alpha}
-
W_tN_t^{(d)}
-
R_tK_{t-1}^{(d)}
+
(1-\delta)K_{t-1}^{(d)}
\right\}
\]

\[
\Rightarrow\qquad
\begin{aligned}
\alpha Z_t(K_{t-1}^{(d)})^{\alpha-1}(N_t^{(d)})^{1-\alpha}+(1-\delta)
&=R_t,
&& \text{(3)}\\[0.6em]
(1-\alpha) Z_t(K_{t-1}^{(d)})^{\alpha}(N_t^{(d)})^{-\alpha}
&=W_t,
&& \text{(4)}
\end{aligned}
\]

System of equations is exactly the same as SP's problem. (isomorphic)

Here, $R_t$ is no longer rental rate, but the gross return factor: rental + resale value.

Define the usual rental rate as
\[
r_t=R_t-(1-\delta)
\]

\[
(3)\Rightarrow MPK=r_t
\]

\[
\Pi_t
=
Z_tK^\alpha N^{1-\alpha}
+(1-\delta)K
-RK-WN
=
Z_tK^\alpha N^{1-\alpha}
-\bigl(R-(1-\delta)\bigr)K-WN
\]
\[
=
Z_tK^\alpha N^{1-\alpha}-rK-WN
\]

\[
(1)\Rightarrow\qquad
C_t^{-\sigma}
=
\beta E_t\left[C_{t+1}^{-\sigma}(r_{t+1}+1-\delta)\right]
\]

\[
(2)\Rightarrow\qquad
C_t+K_t=W_tN_t+r_tK_{t-1}+(1-\delta)K_{t-1}
\]

\[
\Rightarrow\qquad \text{Setting is totally isomorphic.}
\]
\subsection*{Supplementary notes}

\paragraph{Q1. Technology shock: among $C$, $I$, and $K$, which one rises more?}

Under a positive technology shock, $K_t$ does \emph{not} move on impact, because it is predetermined.  
Thus, on impact, both $C_t$ and $I_t$ rise, while $K_t$ remains unchanged.

The more subtle question is whether $C_t$ or $I_t$ rises more. In the model, this is not a pure theoretical certainty independent of parameters: it depends on the parameter values. 

Empirically, however, the data typically show that consumption rises by less, while investment is more volatile.Later in the course, this can be seen more clearly from the impulse response functions.

\paragraph{Q2. What is the difference between the two definitions of return?}

In the baseline setup, $R$ denotes the \emph{capital rent}, namely the rental price.

Later, in Q2, we define
\[
r = R - (1-\delta).
\]
There, $R$ is reinterpreted as the \emph{gross return on capital}, i.e.\ the total return from holding one unit of capital across periods.

Under this notation:
\[
R = r + (1-\delta).
\]

Hence:
\begin{itemize}
    \item $r$ corresponds to the baseline $R$, namely the rental price ;
    \item $R$ in Q2 is the gross return on capital;
    \item the gross return consists of two parts:
    \[
    \text{gross return} = \text{rental price} + \text{residual value of capital}.
    \]
\end{itemize}

\newpage
\subsection*{5. Solve the RBC model}

Find the Steady State (S.S.) of the economy, drop time subscript:
\[
\bar C=\bar Z\bar K^\alpha +(1-\delta)\bar K-\bar K
\qquad \cdots \quad \text{Resource constraint}
\]

\[
1=\beta\left[\alpha \bar Z\bar K^{\alpha-1}+(1-\delta)\right]
\qquad \cdots \quad \text{Euler equation}
\]

Two variables, $(\bar K,\bar C)$ given parameters $(\bar Z,\beta,\alpha,\delta)$
\[
C_t=C_{t+1}=\bar C,\qquad
K_t=k_{t+1}=\bar K,\qquad
Z_t=Z_{t+1}=\bar Z,\qquad
N_t=1,\ \forall t
\]

\[
\bar K
=
\left[
\frac{\alpha \bar Z}{\frac{1}{\beta}-(1-\delta)}
\right]^{\frac{1}{1-\alpha}},
\qquad
\bar C=\bar Z\bar K^\alpha-\delta \bar K,
\qquad
\bar Y=\bar Z\bar K^\alpha
\]

\[
\frac{\bar C}{\bar K}
=
\bar Z\bar K^{\alpha-1}-\delta
=
\frac{\bar Y}{\bar K}-\delta
=
\frac{\frac{1}{\beta}-(1-\delta)}{\alpha}-\delta
=
\frac{\frac{1}{\beta}-1+(1-\alpha)\delta}{\alpha}
\]

\subsection*{6. Log-linearization}

\textbf{(i) Resource constraint:}
\[
K_t=Z_tK_{t-1}^{\alpha}+(1-\delta)K_{t-1}-C_t
\qquad \Rightarrow \qquad
\bar K=\bar Z\bar K^\alpha +(1-\delta)\bar K-\bar C
\]

\[
\bar K(1+\hat k_t)
=
\bar Z(1+\hat z_t)\bar K^\alpha (1+\hat k_{t-1})^\alpha
+
(1-\delta)\bar K(1+\hat k_{t-1})
-\bar C(1+\hat c_t)
\]

Using
\[
(1+\hat k_{t-1})^\alpha \approx 1+\alpha \hat k_{t-1}
\]
we get
\[
\bar K+\bar K\hat k_t
=
\bar Z\bar K^\alpha
\bigl(1+\hat z_t+\alpha \hat k_{t-1}\bigr)
+
(1-\delta)\bar K+(1-\delta)\bar K\hat k_{t-1}
-\bar C-\bar C\hat c_t
\]

Subtract steady state:
\[
\bar K\hat k_t
=
\alpha \bar Z\bar K^\alpha \hat k_{t-1}
+
\bar Z\bar K^\alpha \hat z_t
+
(1-\delta)\bar K\hat k_{t-1}
-
\bar C\hat c_t
\]

\[
\Rightarrow\qquad
\hat k_t
=
\left(\alpha \frac{\bar Y}{\bar K}+1-\delta\right)\hat k_{t-1}
+
\frac{\bar Y}{\bar K}\hat z_t
-
\frac{\bar C}{\bar K}\hat c_t
\]

\[
\textcolor{black}{\ast\ \text{next period change in } k \text{ equals }}
\]
\[
\textcolor{black}{1)\ \text{change in } K_{t-1}\text{, that result in change in output }(MPK=\alpha \bar Y/\bar K)\ \&\ \text{depreciated capital }1-\delta}
\]
\[
\textcolor{black}{2)\ \text{change in tech } z \text{ that result in change in output}}
\]
\[
\textcolor{black}{3)\ \text{change in consumption } C \text{ that result in change in investment.}}
\]

\textbf{(ii) Euler equation:}
\[
1
=
\beta E_t\left[
\left(\frac{C_t}{C_{t+1}}\right)^\sigma
\bigl(1-\delta+\alpha Z_{t+1}K_t^{\alpha-1}\bigr)
\right]
\]

Define
\[
R_{t+1}=1-\delta+\alpha Z_{t+1}K_t^{\alpha-1}
\qquad \Rightarrow \qquad
\bar R=\frac{1}{\beta}
\]

from
\[
1=\beta E_t\left[\left(\frac{C_t}{C_{t+1}}\right)^\sigma R_{t+1}\right]
\]
we get
\[
1=E_t\left[\left(\frac{C_t}{C_{t+1}}\right)^\sigma \frac{R_{t+1}}{\bar R}\right]
\]

Let
\[
\frac{C_t}{C_{t+1}}
=
\frac{\bar C e^{\hat c_t}}{\bar C e^{\hat c_{t+1}}}
=
e^{\hat c_t-\hat c_{t+1}},
\qquad
\frac{R_{t+1}}{\bar R}=e^{\hat r_{t+1}}
\]

Thus
\[
1
=
E_t\left[
e^{\sigma(\hat c_t-\hat c_{t+1})+\hat r_{t+1}}
\right]
\]

use $e^x\approx 1+x$ for small $x$:
\[
1
\approx
E_t\left[
1+\sigma(\hat c_t-\hat c_{t+1})+\hat r_{t+1}
\right]
\]

\[
\Rightarrow\qquad
\sigma \hat c_t
=
\sigma E_t\hat c_{t+1}
-
E_t\hat r_{t+1}
\tag{1}
\]

Now derive $\hat r_{t+1}$.
Let
\[
MPK_{t+1}=\alpha Z_{t+1}K_t^{\alpha-1}
\qquad \Rightarrow \qquad
\overline{MPK}=\alpha \bar Z\bar K^{\alpha-1}
\]
\[
\bar R = 1-\delta+\overline{MPK}
\]
Now approximate $MPK_{t+1}$ around steady state:
\[
\alpha Z_{t+1}K_t^{\alpha-1}
\approx
\alpha \bar Z(1+\hat z_{t+1})\bar K^{\alpha-1}(1+\hat k_t)^{\alpha-1}
\]
\[
=
\overline{MPK}
\bigl[1+\hat z_{t+1}+(\alpha-1)\hat k_t\bigr]
\]
\[
where (1+x)^a \approx 1+ax\qquad \Rightarrow \qquad
(1+\hat{k}_t)^{\alpha-1} \approx 1+(\alpha-1)\hat{k}_t
\]

\newpage

\[
R_{t+1}
\approx
(1-\delta)+\overline{MPK}
+
\overline{MPK}\bigl[\hat z_{t+1}+(\alpha-1)\hat k_t\bigr]
=
\bar R+\overline{MPK}\bigl[\hat z_{t+1}+(\alpha-1)\hat k_t\bigr]
\]

\[
\hat r_{t+1}
=
\log\left(\frac{R_{t+1}}{\bar R}\right)
\approx
\frac{R_{t+1}-\bar R}{\bar R}
=
\frac{\overline{MPK}}{\bar R}
\bigl[\hat z_{t+1}+(\alpha-1)\hat k_t\bigr]
\]

\[
\frac{\overline{MPK}}{\bar R}
=
\frac{\bar R-(1-\delta)}{\bar R}
=
1-\frac{1-\delta}{\bar R}
=
1-(1-\delta)\beta
\]

Let
\[
\tilde \delta \equiv 1-(1-\delta)\beta
\]
share of gross return that comes from new output rather than keeping old capital.

Then
\[
\hat r_{t+1}
=
\tilde \delta
\bigl[\hat z_{t+1}-(1-\alpha)\hat k_t\bigr]
\tag{2}
\]

Plug (2) into (1):
\[
\sigma \hat c_t
=
\sigma E_t\hat c_{t+1}
-
\tilde \delta E_t\hat z_{t+1}
+
\tilde \delta(1-\alpha)\hat k_t
\]

\[
E_t(\hat c_{t+1}-\hat c_t)
=
\frac{\tilde\delta}{\sigma}
\underbrace{\Bigl(E_t\hat z_{t+1}-(1-\alpha)\hat k_t\Bigr)}_{\textcolor{red}{\text{log deviation of next period's MPK}}}
\]
\[
\textcolor{red}{\text{Expected growth in consumption is proportional to expected returns.}}
\]
\[
\textcolor{red}{
E_t(\hat c_{t+1}-\hat c_t)\propto E_t\widehat{MPK}_{t+1}
}
\]
\[
\textcolor{red}{
MPK\uparrow \Rightarrow \text{want more consumption tomorrow } \Rightarrow \text{save/invest more today } \Rightarrow C \text{ grows}
}
\]
\[
\textcolor{red}{
\sigma\uparrow \Rightarrow \text{stronger smoothing, IES}\downarrow \text{ / more curvature } \Rightarrow \text{households hate changing } C \text{ across time}
}
\]
\[
\textcolor{red}{
\tilde\delta\uparrow \Rightarrow \text{returns move more with MPK}\uparrow \Rightarrow C \text{ responds by more}
}
\]

\textbf{(iii) Technology process:}
\[
\log Z_{t+1}
=
(1-\psi)\log \bar Z+\psi \log Z_t+\varepsilon_t
\]

\[
\log(\bar Z e^{\hat z_{t+1}})
=
(1-\psi)\log \bar Z+\psi \log(\bar Z e^{\hat z_t})+\varepsilon_t
\]

\[
\Rightarrow\qquad
\hat z_{t+1}=\psi \hat z_t+\varepsilon_t
\]
the deviation of technology from S.S. is an AR(1) process

The log-linearized system of equation:
\[
\hat k_t
=
\left(\alpha \frac{\bar Y}{\bar K}+1-\delta\right)\hat k_{t-1}
+
\frac{\bar Y}{\bar K}\hat z_t
-
\frac{\bar C}{\bar K}\hat c_t
\]

\[
\alpha \frac{\bar Y}{\bar K}+1-\delta
=
\overline{MPK}+1-\delta=\bar R=\frac{1}{\beta}
\]

\[
\frac{\bar Y}{\bar K}
=
\frac{\overline{MPK}}{\alpha}
=
\frac{\bar R-(1-\delta)}{\alpha}
=
\frac{\frac{1}{\beta}-(1-\delta)}{\alpha}
=
\frac{\tilde \delta}{\alpha \beta}
\]

\[
\frac{\bar C}{\bar K}
=
\bar Z\bar K^{\alpha-1}-\delta
=
\frac{\overline{MPK}}{\alpha}-\delta
=
\frac{\tilde \delta}{\alpha \beta}-\delta
\]

\[
\Rightarrow\qquad
\begin{cases}
\hat k_t=\dfrac{1}{\beta}\hat k_{t-1}+\dfrac{\tilde\delta}{\alpha\beta}\hat z_t-\left(\dfrac{\tilde\delta}{\alpha\beta}-\delta\right)\hat c_t\\[1.1em]
\sigma E_t\hat c_{t+1}+\tilde\delta(1-\alpha)\hat k_t-\tilde\delta E_t\hat z_{t+1}=\sigma \hat c_t\\[0.8em]
\hat z_{t+1}=\psi \hat z_t+\varepsilon_t
\end{cases}
\]
\[
\textcolor{red}{\ast\ \text{Three equations, three variables } (\hat k_t,\hat c_{t+1},\hat z_{t+1})}
\]

\newpage

\subsection*{7. Solve the system in terms of exogenous parameters \& pre-determined variables}

Pre-determined: $(\hat z_t,\hat k_{t-1})$

Write in matrix form:
\[
\begin{pmatrix}
1 & 0 & 0\\
\tilde\delta(1-\alpha) & \sigma & -\tilde\delta\\
0 & 0 & 1
\end{pmatrix}
\begin{pmatrix}
\hat k_t\\
E_t\hat c_{t+1}\\
E_t\hat z_{t+1}
\end{pmatrix}
=
\begin{pmatrix}
\frac{1}{\beta} & \delta-\frac{\tilde\delta}{\alpha\beta} & \frac{\tilde\delta}{\alpha\beta}\\
0 & \sigma & 0\\
0 & 0 & \psi
\end{pmatrix}
\begin{pmatrix}
\hat k_{t-1}\\
\hat c_t\\
\hat z_t
\end{pmatrix}
\]

Because
\[
\hat z_{t+1}=\psi \hat z_t+\varepsilon_t
\]
is a stationary process, the dynamic property of system is pinned down by the deterministic part:
\[
\textcolor{red}{\hat z_t \text{ is either 0 or converges to 0 since } 0<\psi<1}
\]

\[
M\,E_t[x_t]=N x_{t-1}
\]
where
\[
x_{t-1}=\begin{bmatrix}\hat k_{t-1} & \hat c_t\end{bmatrix}'
\]

\[
M=
\begin{pmatrix}
1 & 0\\
\tilde\delta(1-\alpha) & \sigma
\end{pmatrix},
\qquad
N=
\begin{pmatrix}
\frac{1}{\beta} & \delta-\frac{\tilde\delta}{\alpha\beta}\\
0 & \sigma
\end{pmatrix}
\]

\[
\det(M)=\sigma>0
\qquad \Rightarrow \qquad
M^{-1}\text{ exists}
\]

Assume
\[
M^{-1}=
\begin{pmatrix}
b_{11} & b_{12}\\
b_{21} & b_{22}
\end{pmatrix}
\]

Impose $MM^{-1}=I$:
\[
\begin{pmatrix}
1 & 0\\
\tilde\delta(1-\alpha) & \sigma
\end{pmatrix}
\begin{pmatrix}
b_{11} & b_{12}\\
b_{21} & b_{22}
\end{pmatrix}
=
\begin{pmatrix}
1 & 0\\
0 & 1
\end{pmatrix}
\]

\[
\Rightarrow\qquad
b_{11}=1,\quad
b_{12}=0,\quad
b_{22}=\frac{1}{\sigma},\quad
b_{21}=-\frac{\tilde\delta(1-\alpha)}{\sigma}
\]

\[
\Rightarrow\qquad
M^{-1}
=
\begin{pmatrix}
1 & 0\\
-\dfrac{\tilde\delta(1-\alpha)}{\sigma} & \dfrac{1}{\sigma}
\end{pmatrix}
\]

\[
\Rightarrow\qquad
M^{-1}N
=
\begin{pmatrix}
1 & 0\\
-\dfrac{\tilde\delta(1-\alpha)}{\sigma} & \dfrac{1}{\sigma}
\end{pmatrix}
\begin{pmatrix}
\frac{1}{\beta} & \delta-\frac{\tilde\delta}{\alpha\beta}\\
0 & \sigma
\end{pmatrix}
\]

\[
=
\begin{pmatrix}
\dfrac{1}{\beta} & \delta-\dfrac{\tilde\delta}{\alpha\beta}\\[0.8em]
-\dfrac{\tilde\delta(1-\alpha)}{\sigma\beta}
&
1+\dfrac{\tilde\delta(1-\alpha)}{\sigma}
\left(
\dfrac{\tilde\delta}{\alpha\beta}-\delta
\right)
\end{pmatrix}
\]

Using
\[
\delta-\frac{\tilde\delta}{\alpha\beta}
=
-\frac{\bar C}{\bar K},
\]
we can write
\[
M^{-1}N
=
\begin{pmatrix}
\dfrac{1}{\beta} & -\dfrac{\bar C}{\bar K}\\[0.8em]
-\dfrac{\tilde\delta(1-\alpha)}{\sigma\beta}
&
1+\dfrac{\tilde\delta(1-\alpha)}{\sigma}\dfrac{\bar C}{\bar K}
\end{pmatrix}
\]

\[
E_t[x_t]=M^{-1}N x_{t-1}
\]
\[
\textcolor{red}{\text{Jacobian matrix}}
\]

Thus, $M^{-1}N$ shows the property of roots.

\newpage

\[
T
=
1+\frac{1}{\beta}
+
\frac{\tilde\delta(1-\alpha)}{\sigma}\frac{\bar C}{\bar K}
\]

\[
D
=
\frac{1}{\beta}
\left[
1+\frac{\tilde\delta(1-\alpha)}{\sigma}\frac{\bar C}{\bar K}
\right]
-
\frac{\tilde\delta(1-\alpha)}{\sigma\beta}\frac{\bar C}{\bar K}
=
\frac{1}{\beta}
\]

\[
p(-1)=1+D+T
=
2+\frac{2}{\beta}
+
\frac{\tilde\delta(1-\alpha)}{\sigma}\frac{\bar C}{\bar K}
>0
\]

\[
p(1)=1+D-T
=
-\frac{\tilde\delta(1-\alpha)}{\sigma}\frac{\bar C}{\bar K}
<0
\]

\[
\Rightarrow\qquad
\begin{cases}
p(-1)>0\\
p(1)<0
\end{cases}
\]

\[
\Rightarrow\qquad
1 \text{ stable root in }(-1,1),\qquad
1 \text{ unstable root in }(1,+\infty)
\]

Since we have one predetermined variable $\hat k_{t-1}$ and one jump variable $\hat c_t$,
\[
\Rightarrow\qquad
\text{the system is a saddle path equilibrium}
\]
\[
\text{(from Blanchard and Kahn conditions)}
\]

\subsection*{8. What is Jacobian Matrix?}

Partial derivative that tells how a change in today's state affects tomorrow's variables.

eg:
\[
x_t=f(x_{t-1},z_t)
\]

\[
J=\frac{\partial f}{\partial x}\Big|_{(\bar x,\bar z)}
\]
(Notes transcribed by CMT.)

\ThisNoteEnd
